\documentclass[11pt]{article}

\usepackage[a4paper,margin=0.78in]{geometry}
\usepackage[T1]{fontenc}
\usepackage[utf8]{inputenc}
\usepackage{lmodern}
\usepackage{amsmath,amssymb,amsthm,mathtools}
\usepackage{enumitem,booktabs,array,tabularx,microtype,xcolor}
\usepackage[numbers,sort&compress]{natbib}
\usepackage[colorlinks=true,linkcolor=blue!65!black,
  citecolor=blue!65!black,urlcolor=blue!70!black]{hyperref}
\usepackage[nameinlink,noabbrev]{cleveref}
\setlength{\emergencystretch}{3em}

\hypersetup{
  pdftitle={MVP8 after the Source-Corpus and Fixed-Atom Selector Audits},
  pdfauthor={Liang Wang},
  pdfsubject={A source-locked route decision with scoped structural and phase stops},
  pdfkeywords={MVP8, occurrence edge, active support, fixed atom, route classifier}
}

\newtheorem{theorem}{Theorem}[section]
\newtheorem{proposition}[theorem]{Proposition}
\newtheorem{corollary}[theorem]{Corollary}
\theoremstyle{definition}
\newtheorem{definition}[theorem]{Definition}
\theoremstyle{remark}
\newtheorem{remark}[theorem]{Remark}

\newcommand{\NT}{\textnormal{\textsc{not-testable}}}
\newcommand{\STOP}{\textnormal{\textsc{stop-scoped}}}
\newcommand{\OPEN}{\textnormal{\textsc{open}}}
\newcommand{\Lzero}{\mathrm{L0}}
\newcommand{\Lone}{\mathrm{L1}}
\newcommand{\Ltwo}{\mathrm{L2}}
\newcolumntype{Y}{>{\raggedright\arraybackslash}X}
\newcolumntype{P}[1]{>{\raggedright\arraybackslash}p{#1}}

\title{\textbf{MVP8 after the Source-Corpus and Fixed-Atom}\\
\textbf{Selector Audits: Scoped Stops, Preserved Roots,}\\
\textbf{and the Pointwise Arithmetic Frontier}}
\author{Liang Wang\\
School of Artificial Intelligence and Automation,\\
Huazhong University of Science and Technology, Wuhan 430074, P.R. China\\
\texttt{wangliang.f@gmail.com}}
\date{July 2026}

\begin{document}
\maketitle

\begin{abstract}
TPC-173--181 test two interfaces left open by MVP7.  The structural
side freezes every TPC-133--172 manuscript, maps every
occurrence-adjacent claim, and finds no theorem statement satisfying
the complete source-locked local actual-occurrence edge contract.
Consequently the largest defensible local family in that declared
corpus is empty.  Its gluing coverage is \(0\) of \(2988\) archived
cuts.  Empty eligibility does not certify actual active support, and
the five-field archive address is not a canonical physical
representative.  The extraction cell for this corpus is stopped, but
the occurrence-augmented architecture remains open and
not-testable.

The phase side source-locks the fact \(h_0=2\), but finds no named
physical atom, phase value, or production packet-coordinate row.
Moreover, an almost-every-phase theorem cannot select a prescribed
singleton: a full-measure good set may omit that point.  Thus the
uncontrolled metric-to-atom method is stopped in its exact scope.
The source-backed selector remains not-testable, and the two direct
pointwise fixed-atom targets remain open.

MVP8 imports these results without promoting their quantifiers.  Its
verdict is \(\NT\), with first missing object
\[
\mathsf{H1.source\_backed\_local\_occurrence\_edge\_family}.
\]
The seven-node structural/physical blocker antichain is unchanged.
The TPC-170 phase-metric powers remain ineligible for a named-atom
endpoint charge.  No program-positive \(\Ltwo\), strict \(1/400\),
prime-pair lower bound, or twin-prime theorem is proved.
\end{abstract}

\section{Dynamic imported state}

MVP7 ended with three structural roots and four physical registry
roots \citep{WangTPC172}.  The present classifier does not assume that
the planned TPC-173--181 route succeeded.  It reads the canonical
artifacts of those papers and rejects any mismatch in family
cardinality, coverage, support status, representation status, phase
identity, selector status, or pointwise-route status.  The resulting
eight source hashes are integrity locks only.

Let
\[
\begin{split}
\mathcal B_{\mathrm{H1}}=\{&
\mathsf{source\_backed\_local\_occurrence\_edge\_family},\\[-2pt]
&\mathsf{actual\_active\_support\_certificate},\\[-2pt]
&\mathsf{canonical\_minimal\_representation\_certificate}\}.
\end{split}
\tag{1}\label{eq:h1-roots}
\]
The physical roots are
\[
\mathcal B_{\mathrm{H9}}=
\{\mathsf{literal\_weight},\mathsf{phase\_cell},
\mathsf{endpoint},\mathsf{normalization}\}_{\rm registry}.
\tag{2}\label{eq:h9-roots}
\]
All seven nodes have status \(\NT\).  Their roles differ: the first
three are structural certificates and the last four are literal data
registries with decay axis \(\mathsf{NONE}\).

\section{The scoped structural stop}

TPC-173 freezes the contiguous theorem corpus
\[
\mathcal C_{133:172}
=\{\text{the unique \texttt{main.tex} of TPC-}n:
133\leq n\leq 172\}.
\tag{3}\label{eq:corpus}
\]
Its forty files are partitioned into thirty mapped but disqualified
sources and ten reviewed files with no candidate; no file is left
unmapped.  A qualifying claim must have resolving theorem and formula
locators, a nonempty derivation, an actual local-occurrence
conclusion, the five-field production cut address, an exact edge
weight, fixed-\(h_0=2\) lineage, and physical-normalization lineage.

\begin{proposition}[Declared-corpus maximality]\label{prop:empty}
Relative to \(\mathcal C_{133:172}\), the largest source-locked local
actual-occurrence edge family has cardinality zero.
\end{proposition}

\begin{proof}
TPC-173 maps every candidate and obtains zero qualifying claims.
TPC-174 makes the required witness contract explicit, while keeping
its two-edge nonvacuity fixture synthetic.  TPC-175 requires every
production row to import a qualifying TPC-173 claim.  Hence any
nonempty admissible family would contain a row whose source claim is
in the empty qualifying set, a contradiction
\citep{WangTPC173,WangTPC174,WangTPC175}.
\end{proof}

The conclusion of \cref{prop:empty} is closed-world.  It is not
\[
\text{``no local actual-occurrence edge exists in mathematics.''}
\tag{4}\label{eq:not-global}
\]
A new source corpus or a genuinely new theorem can define a new
extraction cell.

TPC-176 applies the formal gluing interface only to the proved local
family.  The exact coverage ledger is
\[
\#\mathrm{proved\ local\ edges}=0,\qquad
\#\mathrm{covered\ cuts}=0,\qquad
\#\mathrm{unmatched\ cuts}=2988.
\tag{5}\label{eq:coverage}
\]
There are no duplicates because there are no eligible rows.  This
does not satisfy the nonempty local-totality hypotheses of the
TPC-165 gluing theorem \citep{WangTPC165,WangTPC176}.

\section{Vacuity is not a certificate}

The other two roots in \eqref{eq:h1-roots} are independent of local
edge existence.  TPC-177 and TPC-178 therefore do not close them by
vacuity.

\begin{proposition}[Support firewall]\label{prop:support}
Zero sourced carriers do not prove an actual active-support
certificate.  In particular, the statement
\[
\forall o\in\varnothing,\quad c_{\rm phys}(o)\ne0
\tag{6}\label{eq:vacuous}
\]
is vacuously true but does not identify the required production
carrier or its nonzero literal coefficient.
\end{proposition}

\begin{proof}
The target certificate is existentially anchored: it must source-lock
the actual carrier and then prove nonvanishing of its literal
coefficient.  Universal truth on an empty test domain supplies neither
object.  The separate H9 literal-weight registry also remains missing;
it cannot be inferred from a support predicate \citep{WangTPC177}.
\end{proof}

Likewise, TPC-164's unique minimum separating key
\[
(\ell,k,\mathrm{native\_d},j_L,j_K)
\tag{7}\label{eq:archive-key}
\]
is an archive address.  TPC-178 finds no actual carrier on which to
test a physical quotient or minimal parent.  Thus address injectivity
does not imply canonical physical representation
\citep{WangTPC164,WangTPC178}.

TPC-179 integrates these facts.  The corpus extraction cell is
\(\STOP\), while the global occurrence-augmented architecture is
\(\OPEN/\NT\).  The three-root antichain \eqref{eq:h1-roots} is
unchanged \citep{WangTPC179}.

\section{The production phase registry}

TPC-180 separates one positive data fact from the missing phase
registry:
\[
h_0=2\quad\text{is source-backed},\qquad
\alpha_{\rm phys},\quad
\operatorname{loc}(\alpha_{\rm phys}),\
\mathcal P_{\rm prod}
\quad\text{are absent}.
\tag{8}\label{eq:phase-census}
\]
Here \(\mathcal P_{\rm prod}\) is the exact source-locked list of
production packet coordinates.  The census has no named atom ID, no
phase value modulo one, and no packet-coordinate row.  Consequently
\[
\mathsf{H9.phase\_cell\_registry}=\NT,\qquad
\mathsf{decay\_axis}=\mathsf{NONE}.
\tag{9}\label{eq:phase-registry}
\]
Neither the value \(h_0=2\) nor a future registry would manufacture a
cancellation exponent \citep{WangTPC180}.

The strongest arithmetic import remains the TPC-170 theorem:
for Lebesgue-almost every fixed phase, eventually every declared
packet and every prefix on the prescribed schedule obeys the
phase-metric bound, with powers \(X^{-\delta}\) for every
\(\delta<1/4\) in the polylogarithmic regime
\citep{WangTPC170}.  The exceptional null set depends on the entire
schedule.

\section{The selector obstruction}

\begin{theorem}[Singleton nonimplication]\label{thm:singleton}
Let \(\alpha_\star\) be any prescribed point of the phase torus.
The assertion that a measurable good set has full Lebesgue measure
does not imply that it contains \(\alpha_\star\).
\end{theorem}

\begin{proof}
The set
\[
\mathbb T\setminus\{\alpha_\star\}
\tag{10}\label{eq:singleton}
\]
has full Lebesgue measure and omits \(\alpha_\star\).  Equivalently, a
measure-zero exceptional set may be precisely the prescribed
singleton.  Applying this logical model to the TPC-170
Borel--Cantelli conclusion shows that the metric theorem alone does
not prove avoidance of its schedule-dependent limsup exceptional set
at a named physical atom.
\end{proof}

TPC-181 therefore records
\[
\begin{array}{rcl}
\mathsf{phase\_metric\_uncontrolled\_atomic}
&=&\STOP,\\
\mathsf{phase\_metric\_source\_backed}
&=&\NT.
\end{array}
\tag{11}\label{eq:selector}
\]
The first line is a scoped logical obstruction.  The second can be
reopened by a source-locked named phase, its exact production schedule,
and a theorem that this point avoids the corresponding exceptional
set.  A direct pointwise fixed-phase theorem bypasses the metric
selector entirely \citep{WangTPC181}.

\section{Typed routes and endpoint ledger}

Architecture and arithmetic methods remain disjoint.  The current
route cells are summarized in \cref{tab:routes}.

\begin{table}[ht]
\centering
\small
\caption{MVP8 route cells.  A scoped cell stop does not stop its
parent architecture or an independent theorem method.}
\label{tab:routes}
\begin{tabularx}{\textwidth}{P{0.35\textwidth}P{0.20\textwidth}Y}
\toprule
Cell & State & Exact meaning\\
\midrule
Occurrence-augmented architecture
  & \(\OPEN/\NT\)
  & Still lacks all three structural roots.\\
TPC-133--172 edge extraction
  & \(\STOP\)
  & No qualifying source claim in that frozen corpus.\\
Uncontrolled metric-to-atom promotion
  & \(\STOP\)
  & Almost-every phase does not select a prescribed singleton.\\
Source-backed metric selector
  & \(\OPEN/\NT\)
  & Named atom, production schedule, and avoidance theorem are missing.\\
Bad-endpoint pointwise fixed atom
  & \(\OPEN\)
  & Parent-ready O161 arithmetic target.\\
Direct additive-twist fixed atom
  & \(\OPEN\)
  & Parent-ready O161 arithmetic target.\\
\bottomrule
\end{tabularx}
\end{table}

Endpoint Ledger V5 keeps the metric and named-atom powers separate:
\[
\sigma_{\rm metric}=\{\delta:\delta<1/4\},
\qquad
\sigma_{\rm named}=0,
\qquad
\sigma_{\rm req}=\frac1{400}.
\tag{12}\label{eq:ledger}
\]
The first entry has phase quantifier
\(\mathsf{LEBESGUE\_AE\_FIXED\_PHASE}\), so it is not eligible for the
second entry.  Literal coefficients, atomic normalization, actual
support, canonical/minimal representation, and strict endpoint slack
remain not-testable.  Therefore the strict \(1/400\) charge is unpaid.

All six quantifier axes remain unpromoted:
\[
\begin{array}{c|c|c}
\text{axis}&\text{proved}&\text{required}\\
\hline
\text{carrier}&\text{explicit packet corridor}&\text{actual fixed-}h_0\text{ packet}\\
\text{phase}&\text{Lebesgue-a.e. fixed phase}&\text{named fixed atom}\\
\text{endpoint}&\text{all shell prefixes}&\text{deterministic all prefixes}\\
\text{scale}&\text{eventual prescribed schedule}&\text{deterministic all scales}\\
\text{decay}&\text{phase-metric }X\text{-power}&\text{fixed-atom }X\text{-power}\\
\text{support}&\text{actual core}&\text{actual active support}
\end{array}
\tag{13}\label{eq:projection}
\]

\section{MVP8 decision}

\begin{theorem}[MVP8]\label{thm:mvp8}
For the source-locked TPC-172 and TPC-175--181 artifacts, the current
route verdict is
\[
\boxed{\mathsf{NOT\_TESTABLE}}.
\tag{14}\label{eq:mvp8}
\]
The first missing object is
\(\mathsf{H1.source\_backed\_local\_occurrence\_edge\_family}\), and
the minimal not-testable antichain is exactly
\(\mathcal B_{\mathrm{H1}}\cup\mathcal B_{\mathrm{H9}}\).
The two O161 pointwise fixed-atom targets remain open and
parent-ready.
\end{theorem}

\begin{proof}
\Cref{prop:empty,prop:support} and the TPC-178 representation audit
leave all three H1 roots not-testable.  Equations
\eqref{eq:phase-registry} and \eqref{eq:selector} leave every H9 root
missing and stop only one arithmetic method cell.  No architecture
infeasibility theorem or complete route-universe theorem is imported.
The ledger \eqref{eq:ledger} has no named-atom power and no strict
endpoint slack.  The classifier therefore selects \(\NT\), preserving
the declared first-root order and both disjoint pointwise targets.
\end{proof}

The next structural object is not another hash or schema: it is a new
explicit source corpus or a new theorem-backed local occurrence edge.
The next phase object is a source-locked named atom together with a
schedule-specific avoidance theorem; absent that bridge, work returns
to the two pointwise O161 targets.

\section{Claim firewall}

The batch contributes \(\Lzero\) executable contracts and diagnostics,
and two \(\Lone\) scoped obstructions: declared-corpus source
exhaustion and the metric-to-singleton nonimplication.  It contributes
no new production occurrence family, active-support certificate,
canonical physical representation, named fixed-atom theorem, or
program-positive \(\Ltwo\) estimate.

In particular, this paper does not prove RH, a strict \(1/400\)
endpoint gain, \(B_{2,\delta}(X)=o(X)\), a prime-pair lower bound, the
twin-prime conjecture, or infinitude of twin primes.

\bibliographystyle{plainnat}
\bibliography{references}

\end{document}
